% Nouveau EGA V (cote 25): archive cote 25, site page 176 (PDF page 177).
% typescript. Draft compared: none.
% AI-transcribed by Claude Opus 5.5 (claude-opus-5-5, adaptive thinking, effort high), Message Batches API, batch ega25b (msgbatch_01A7n5YuXubUPNK1zsepi8RG), 2026-11-13; not yet checked by eye.
\scanpage{25/176}
\mspage{\hand{8}}

coéfficients entiers en les $a_i,b_{ij}$ . Alors le ``contenu'' de ce polynôme
est égal à 2 (i.e. le pgcd de ses coéfficients est 2).
\var{\inserted[typed]{multiple de 2.}}{multiple de 2.}{insertion}{ega25_25_176-v0001}
\var{\xout{Cela}}{}{typewriter strike-out}{ega25_25_176-v0002} Prouvons d'abord que ce contenu est \var{\xout{de la forme 2\var{\deleted{$^r$}}{}{deletion}{ega25_25_176-v0004}\illegible{1}}}{}{typewriter strike-out}{ega25_25_176-v0003}
En effet, cela signifie aussi que \var{\deleted{\var{\xout{\illegible{1}}}{}{typewriter strike-out}{ega25_25_176-v0006}}}{}{deletion}{ega25_25_176-v0005} lorsqu'on spécialise ce polynôme
en \var{\xout{la}}{}{typewriter strike-out}{ega25_25_176-v0007}un corps $Z/\var{\over{p}{2}}{2}{writing over}{ega25_25_176-v0008}Z$\hand{)} \var{\deleted{avec $p\neq 2$}}{}{deletion}{ega25_25_176-v0009}, il \var{\xout{n'}}{}{typewriter strike-out}{ega25_25_176-v0010}est pas identiquement nul, ce qui \var{\xout{xxx}}{}{typewriter strike-out}{ega25_25_176-v0011}
résulte du fait qu'une forme quadratique de degré imp\var{\over{q}{a}}{a}{writing over}{ega25_25_176-v0012}ir sur un corps de
car 2 est tjrs dégénérée, i.e. à un discriminant nul. Pour prouver
que le contenu \var{\inserted{c}}{c}{insertion}{ega25_25_176-v0013} est exactement 2, il suffit de fai\unsure{r}e le calcul du discri-%
minant \var{\inserted{de la forme}}{de la forme}{insertion}{ega25_25_176-v0014} standart de degré n, d'une part elle doit être multiple de c,
d'autre part le calcul donne 2, (car c'est 1 pour le degré pair, et 2 pour
le degré 1 ...).

\var{\xout{\unsure{On en conclut immédiatement}}}{}{typewriter strike-out}{ega25_25_176-v0015} Il y a lieu alors d'introduire le polynôme
\[ \overset{\hand{\sim}}{d}\var{\xout{*}}{}{typewriter strike-out}{ega25_25_176-v0016}(Q) = \tfrac{1}{2}d(Q) \]
en les coéfficients de Q, qui est encore à coéfficients entiers, donc
prend une valeur bien déterminée quand on spécialise les $a_i,b_{ij}$ dans un
anneau A quelconque, i.e. quand on prend une forme quadratique \var{\deleted{\var{\xout{Q}}{}{typewriter strike-out}{ega25_25_176-v0018}}}{}{deletion}{ega25_25_176-v0017}\var{\inserted{q}}{q}{insertion}{ega25_25_176-v0019}de degr\unsure{é}
n à coéfficients dans A. On appela $\overset{\hand{\sim}}{d}\var{\xout{*}}{}{typewriter strike-out}{ega25_25_176-v0020}(Q)$ le``\uline{polynôme} \unsure{\hand{``}}\uline{discrimihant}
\uline{corrigé de la forme quadratique indétérminée Q à n variables}'' , et sa
valeur relativement aux coéfficients d\var{\xout{e}}{}{typewriter strike-out}{ega25_25_176-v0021}'une forme quadratique q à coeff
dans un anneau A quelconque sera appelé le ``\uline{discriminant corrigé de la}
\uline{forme quadratique} q''. Plus généralement, on déduit de façon essentielle
ment triviale de 16.5. l'énoncé suivant:

\uline{Proposition} 16.6. On peut, \var{\inserted{et}}{et}{insertion}{ega25_25_176-v0022} d'une seule manière, associéer à tout\var{\xout{e}}{}{typewriter strike-out}{ega25_25_176-v0023}
préschéma S, muni d'un Module loc libre \var{\xout{d}}{}{typewriter strike-out}{ega25_25_176-v0024}e rang fini \uline{E} et d'une forme \var{\inserted{qu.}}{qu.}{insertion}{ega25_25_176-v0025}
$Q \in \unsure{\Gamma}\,\underline{\mathrm{Sym}}^2(E)$, associer une section $\overset{\hand{\sim}}{d}(Q)$ de dét$(E)^{\otimes 2}$ , de façon à satis-%
faire aux conditions suivantes:

(\var{\xout{1}}{}{typewriter strike-out}{ega25_25_176-v0026}a) Compatibilité avec changement de base\hand{, et fonctorialité pour isomorphismes en \uline{E}.}

(b) Si E est partout de rang pair, on a $\overset{\hand{\sim}}{d}(Q) =$ d(Q) (ou d(Q)
désigne le discriminant de Q).

\mspage{\hand{176}}
